\section{Resumo das contribuições}

As principais contribuições deste trabalho foram:

\begin{itemize}
    \item Desenvolvimento de uma modelagem integrada que combina dinâmica orbital no referencial \gls{eci}, movimento relativo no referencial \gls{lvlh} e dinâmica acoplada entre espaçonave e \gls{mr}.
    
    \item Formulação dinâmica baseada em abordagem lagrangiana, que possibilitou a existência dos torques de reação da espaçonave devido ao movimento do \gls{mr} em seu modo operacional.
    
    \item Implementação computacional e simulação em Matlab e Simulink.
    
    \item Integração de leis de controle \gls{pd} aplicadas ao controle de atitude da espaçonave, à compensação da espaçonave diante aos torques de reação gerados pelo \gls{mr}, ao controle das juntas do manipulador e ao controle do movimento relativo orbital baseado nas equações de \gls{hcw}.
    
    \item Definição e aplicação de critérios numéricos para posição, velocidade relativa, atitude e de inserção da ferramenta, utilizados como hipóteses de simulação para análise do comportamento dinâmico.
\end{itemize}